\section{積分可能性}\label{節:積分可能性}
定義では,リーマン和の極限が存在するときにその値を定積分と呼ぶと決めた.しかし,極限が存在することそのものはまだ保証していない.すべての分割と代表点を調べることはできないから,どのような関数なら積分可能かを先に知っておきたい.本書で主に扱う連続関数については,次の定理を用いる.

\begin{theorem}{連続関数の積分可能性}\label{定理:連続積分可能}
閉区間 $[a,b]$ で連続な関数は,その区間でリーマン積分可能である.
\end{theorem}

この定理の完全な証明は,本書では扱わない.考え方だけ述べておこう.十分短い小区間では,関数の高さの変動も小さくなる.各小区間での最大の高さと最小の高さの差が $\eta$ 以下なら,各小区間で最大の高さを使った和と最小の高さを使った和の差は
\[
 \sum_{k=1}^{n}\eta\Delta x_k=\eta(b-a)
\]
以下になる.どの代表点を選んだリーマン和もこの2つの和の間にあるから,リーマン和どうしの差も $\eta(b-a)$ 以下である.閉区間上の連続性を使うと,\,$\|P\|$ を小さくすることで,この $\eta$ を区間全体で一斉に小さくできる.このような評価をすべての分割について保証する部分も含め,上の定理として採用する.

この定理があれば,連続関数の定積分を計算するときは,計算しやすい分割と代表点を1組選べばよい.極限の存在が保証されているので,どれを選んでも同じ値に近づくからである.\,$[a,b]$ を $n$ 等分して右端を選ぶと,幅は $\frac{b-a}{n}$,第 $k$ 小区間の右端は $a+\frac{k(b-a)}n$ だから,
\[
 \int_a^b f(x)\dd x
 =\lim_{n\to\infty}\frac{b-a}{n}
  \sum_{k=1}^{n}f\left(a+\frac{k(b-a)}n\right)
\]
となる.

\begin{bookpoint}
\textbf{POINT}\par
連続関数なら,\,$n$ 等分して右端 $a+\frac{k(b-a)}n$ の高さで和を作り,\,$\sum$ の公式で計算してから $n\to\infty$ とする.
\end{bookpoint}

\begin{example}{リーマン和}
区間 $[1,3]$ を $n$ 等分し,各小区間の右端を使って $\int_1^3(2x-1)\dd x$ を求めよ.
\end{example}
\paragraphbf{方針}
幅 $\frac2n$,右端 $1+\frac{2k}n$ の和を作り,\,$\sum k$ の公式で計算して $n\to\infty$ とする.

\paragraphbf{解答}
$\int_1^3(2x-1)\dd x=6$

\paragraphbf{解法}
被積分関数は連続なので,定積分は存在する.幅は $\frac2n$,第 $k$ 小区間の右端は $1+\frac{2k}{n}$ である.したがって,リーマン和は
\begin{align*}
 S_n
 &=\frac2n\sum_{k=1}^{n}\left\{2\left(1+\frac{2k}n\right)-1\right\}\\
 &=\frac2n\sum_{k=1}^{n}\left(1+\frac{4k}n\right)\\
 &=\frac2n\left\{n+\frac4n\cdot\frac{n(n+1)}2\right\}\\
 &=2+\frac{4(n+1)}{n}\\
 &=6+\frac4n
\end{align*}
となる.\,$n\to\infty$ とすれば,定積分は $6$ である.

この例では和の公式が使えたが,一般の関数では和を直接計算するのは難しい.後で微分との関係を使い,別の計算法を導く.
